\section{Background}
\label{sec:background}

\subsection{Probabilistic Genotyping Software Systems}
\label{sec:background-pgs}
Forensic DNA testing aims to determine whether a person of interest (e.g., the defendant) could be the source of some or all of the DNA recovered from a crime scene~\cite{butler2024dna}. Since 2009, the U.S. criminal legal system has increasingly relied on \textbf{probabilistic genotyping software (PGS)} to analyze complex DNA evidence~\cite{moss2015admissibility, foley}.\footnote{\citet{moss2015admissibility} describes, ``In 2009, TrueAllele made its first appearance in a United States courtroom in the murder case of Commonwealth v. Foley.''}. PGS are now widely used in criminal cases around the country. For example, in March of 2025, the creators of STRmix---one of the two leading PGS alongside TrueAllele---announced that STRmix is being used by 91 forensic labs across the U.S., by 29 more labs internationally, and in over 690,000 criminal cases worldwide~\cite{strmix2025use}
 

Figure~\ref{fig:pgs-flowchart} depicts how forensic labs integrate PGS in casework. A laboratory analyst first prepares two DNA samples for input into the software: an \emph{evidentiary sample} taken from the crime scene and a \emph{reference profile} from a person of interest. The analyst then provides these samples as inputs to the software alongside subjective parameters, most notably the analyst's best estimate for the ``apparent number of contributors.'' The PGS uses these inputs to produce a \textbf{likelihood ratio (LR)}. This final output is significantly impacted by the analyst's decision-making, such as their estimate for the number of contributors, as well as upstream laboratory equipment~\cite{butler2024dna, paoletti2005empirical,coble2015uncertainty}.\footnote{For more background on likelihood ratios and their use in forensic DNA statistics, see Section 2.5, ``Likelihood Ratios: Introduction to Theory and Application,'' in \citet{butler2024dna}}. 


The LR compares the probability of observing the evidentiary sample under two competing hypotheses: one where the person of interest is a contributor and one where they are not.\footnote{For example, an LR may compare a hypothesis that the person of interest and two unknown individuals contributed to the DNA mixture, against a hypothesis that three unknown individuals who are not the person of interest contributed to the DNA mixture.} LR values greater than 1 are interpreted as providing support for \textbf{inclusion}, while values less than 1 support \textbf{exclusion}. An $LR = 1$ is considered ``uninformative.'' We define a false negative result as an exclusionary LR ($LR < 1$) for a true contributor (i.e., a person of interest whose DNA \textit{is} in the evidentiary sample) and a false positive as an inclusionary LR ($LR > 1$) for a non-contributor (i.e., a person of interest whose DNA is \textit{not} in the evidentiary sample). 


\subsection{Validation of PGS Systems}
\label{sec:background-pgs-validation}
In forensics, validation is the empirical testing of a method prior to its deployment in casework. For PGS, this process consists of two parts: \emph{developmental validation} (conducted by software developers) and \emph{internal validation} (conducted by forensic laboratories seeking to use the software in casework)~\cite{butler2024dna, coble2019probabilistic}. In this work, we focus on internal validation studies, which we conceptualize as \textit{audits}\footnote{In this paper, we use \citet{birhane2024ai}'s definition of an audit as ``any independent assessment of an identified audit target via an evaluation of articulated expectations with the implicit or explicit objective of accountability.'' For a detailed explanation of the terms used in this definition, see Section II.A in \citet{birhane2024ai}.}. 


\subsubsection{Purpose of internal validation studies} 
A forensic lab must complete an internal validation study before using a PGS in casework. A goal of internal validation is to establish the operational range---the range of DNA profile characteristics for which the lab has demonstrated acceptable performance~\cite{butler2024dna}. Dozens of profile characteristics are known to impact the reliability of PGS. Examples of those known to have the greatest impact include the number of contributors to the sample (NoC), the ratio of DNA amounts of contributors to the sample (if the sample is a DNA mixture, i.e., has more than one contributor), and the total amount of DNA in the sample~\cite{butler2024dna}.\footnote{For a list of notable profile characteristics, see Table 4.1 in \citet{butler2024dna}.} 


Internal validation studies and the operational range established using study results can play a critical role in judges' decisions about whether or not the PGS outputs are sufficiently reliable to be presented as evidence at trial~\cite{roth2016machine}. This has been the case in several federal courts. For instance, the courts in U.S. v. Ortiz~\cite{ortiz} and U.S. v. Johnston~\cite{johnston} excluded PGS outputs due to sample characteristics in these cases exceeding the operational bounds established by the forensic laboratories' internal validation studies. Likewise, the court in U.S. v. Lewis~\cite{2020lewis} admitted some evidence and excluded some evidence after similar considerations. An internal validation study's compliance with an established standard can provide additional assurance to a judge that the forensic lab performed a high quality internal validation study and that the operational range is reliable.


\subsection{Standards and Guidelines for PGS Validation}
\label{sec:background-pgs-guidance}
In 2009, the National Academy of Sciences issued a report identifying a wide range of shortcomings in U.S. forensic science research and practice~\cite{national2009strengthening}. Recommendations six to eight of this report emphasize the need to establish forensic science standards, to accredit laboratories adhering to those standards, and to establish quality assurance and quality control practices for individual practitioners in order to promote rigor and reliability in forensic science methods. 

Following the NAS report, standards have been developed for a range of forensic disciplines. However, recent work raises the concern that these efforts have not yielded meaningful improvements and instead threaten to maintain the status quo~\cite{sinha2022radically, morrison2020vacuous}. We build on this work by exploring how standards for internal validation of probabilistic genotyping software may contain similar deficiencies. While several national and international groups have issued guidance for PGS validation~\cite{swgdam2015guidelines, coble2016dna, ukfsr2024guidelines, enfsi2023guidelines, fbi2025qas}, unlike ASB 018, these standards and guidelines are not approved by or registered with the American National Standards Institute (ANSI).

\subsubsection{The AAFS Standards Board (ASB)} 
\label{sec:background-asb}
In 2015, the American Academy of Forensic Sciences established the \textbf{AAFS Standards Board (ASB)}, a standards development organization tasked with developing standards for forensics practice in the U.S.\footnote{For more background on the creation of the ASB, see Appendix~\ref{appendix-a}.} Importantly, as opposed to the guidelines referenced in the previous paragraph, ASB standards are approved by and registered with the ANSI. This means that the ASB develops standards following ANSI requirements for ``openness, balance, lack of dominance, due process, and consensus''~\cite{ansi}. These requirements seek to ensure that the standards development process is ``equitable, accessible, and responsive to the requirements of various stakeholders,'' and that standards developed by the ASB allow organizations to ``maintain autonomy from vested interest groups''~\cite{asb2025about}.


\subsubsection{ASB Standard 018} 
\label{sec:background-asb018}
In 2020, the ASB published the first edition of \textbf{ASB 018, the Standard for Validation of Probabilistic Genotyping Systems}\footnote{ASB 018 is publicly available for download at: \url{https://www.aafs.org/sites/default/files/media/documents/018_Std_e1.pdf}}, which outlines a set of requirements for developmental and internal validations of PGS~\cite{asb018}. 

ASB 018 can influence internal validation study practices in several ways. On one hand, it can provide guidance to labs conducting internal validation studies. On the other hand, it can shape internal validation studies through external scrutiny that pressures labs to ensure internal validation compliance with the standard. The latter currently happens primarily through admissibility hearings in criminal cases, where a judge may interpret a lab's claim of compliance with ASB 018 as providing assurance of validation study quality.\footnote{When PGS is introduced in a criminal case in the U.S., the court may hold an admissibility hearing to determine whether the evidence is sufficiently reliable to be presented in front of the jury. All federal courts and some state courts use the Daubert Standard to conduct this assessment, and one factor considered under the Daubert Standard is \textit{``the existence and maintenance of standards controlling its operation.''}~\cite{wexlaw_daubert} The U.S. v. Anderson court conducted a Daubert hearing for one PGS software, and concluded: ``\textit{the Government has established that [the PGS tool] complies with the relevant standards issued by [SWGDAM] and ANSI/ASB. Accordingly, this factor weighs in favor of admissibility}''~\cite{2023anderson}.} 

In this paper, we focus ASB 018's effectiveness as a tool for assessing and providing assurance of internal validation study quality. Our investigation of ASB 018's effectiveness in providing appropriate assurance about compliant studies builds on critical concerns raised by recent work highlighting that judges frequently hold uncritical perceptions of forensic software~\cite{jin2024beyond} and work raising concerns that standards create ``a danger [...] that a court may not look further than the fact that a standard exists, and be misled into believing that conformity to a vacuous standard is indicative of scientific validity, even though it is not''~\cite[p.207]{morrison2020vacuous}.


